\documentclass[11pt]{article}
\usepackage[T1]{fontenc}
\usepackage{lmodern}
\usepackage[margin=1in]{geometry}
\usepackage{amsmath,amssymb,amsthm,mathtools}
\usepackage[hidelinks]{hyperref}
\usepackage{enumitem}
\usepackage{graphicx}
\emergencystretch=1em

\newtheorem{theorem}{Theorem}
\newtheorem{lemma}[theorem]{Lemma}
\newtheorem{corollary}[theorem]{Corollary}
\newtheorem*{remark}{Remark}
\newcommand{\G}{\Gamma}
\newcommand{\C}{\mathbf C}
\newcommand{\tr}{\operatorname{tr}}
\newcommand{\Tr}{\operatorname{Tr}}
\newcommand{\HS}{\mathrm{HS}}
\newcommand{\one}{\mathbf 1}

\title{The Finite-Compression Determinant Infimum\\Characterizes Amenability}
\author{Run \texttt{fa\_banach\_001}, lane 13}
\date{July 21, 2026}

\begin{document}
\maketitle

\begin{abstract}
Austin asks whether a countable group must be amenable if his finite-subset
determinant lower bound is always sharp for positive definite functions of
the form $\phi=\tau(a^*(\,\cdot\,)a)$ with $a\in\C[\G]$.  We give an
affirmative answer.  The main ingredient is a quantitative form of the
operator Jensen inequality for compressions of $\log$.  Applied to
$a=1+t\sum_{s\in S}s$, it bounds the logarithmic determinant gap below by a
fixed multiple of the normalized $S$-boundary of the compressing set.  Thus
sharpness for these elementary group-ring elements yields the F\o lner
criterion.  The same estimate gives a direct quantitative answer to the
following problem in the source for this family of elements.
\end{abstract}

\section{Source question and answer}

Let $\lambda$ be the left regular representation of a countable group $\G$,
let $\xi=\delta_e$, and let $\tau$ be the regular trace.  For a scalar
positive definite function $\phi$ and finite nonempty $F\subset\G$, put
\[
 \phi[F]=[\phi(g^{-1}h)]_{g,h\in F}.
\]
Austin proves the general lower bound
\begin{equation}\label{eq:source-lb}
 \Delta\phi_{\mathrm{ac}}
 \leq \inf_{\varnothing\ne F\Subset\G}
          (\det\phi[F])^{1/|F|}.
\end{equation}
On PDF page 36 of arXiv:2412.13751v4, Problem 3.7 asks:
\begin{quote}
Suppose that equality holds in \eqref{eq:source-lb} whenever
$\phi=\tau(a^*(\,\cdot\,)a)$ for some $a\in\C[\G]$.
Must $\G$ be amenable?
\end{quote}
See Figure~\ref{fig:problem}.

\begin{theorem}[Full answer to Problem 3.7]\label{thm:main}
Yes.  If equality in \eqref{eq:source-lb} holds for every
$\phi=\tau(a^*(\,\cdot\,)a)$ with $a\in\C[\G]$, then $\G$ is amenable.
In fact, it is enough to assume equality for the positive elements
\[
 a=1+t\sum_{s\in S}s,
 \qquad S=S^{-1}\Subset\G\setminus\{e\},\qquad 0<t<|S|^{-1}.
\]
\end{theorem}

\begin{figure}[ht]
\centering
\includegraphics[width=.96\textwidth]{figures/open_problem_crop.png}
\caption{Equation (3.8) and Problem 3.7 in Austin, arXiv:2412.13751v4,
PDF p.~36.}
\label{fig:problem}
\end{figure}

\section{Proof intuition}

The matrix $\phi[F]$ is a coordinate compression of a positive
right-convolution operator $B$.  Since every diagonal matrix coefficient of
$\log B$ is the von Neumann trace, the logarithm of the ratio between the two
sides of \eqref{eq:source-lb} is exactly the Jensen gap
\[
 |F|^{-1}\bigl(\Tr\log(P_FBP_F)-\Tr(P_F\log B P_F)\bigr).
\]
The resolvent formula for $\log$ and a Schur complement show that this gap
controls the Hilbert--Schmidt norm of the off-diagonal block
$(1-P_F)BP_F$.

Now prescribe a finite symmetric $S$ and choose
$a=1+t\sum_{s\in S}s$.  The associated positive right-convolution operator
is $T=1+tH$, and $B=T^2$.  Positivity of $T$ makes the Sylvester map
$Y\mapsto CY+YA$ uniformly invertible, so leakage of $T^2$ controls leakage
of $H$.  But the columns of $H$ are sums of distinct coordinate vectors;
hence its leakage is exactly
\[
 \|(1-P_F)HP_F\|_{\HS}^2=\sum_{s\in S}|Fs\setminus F|.
\]
Therefore sharpness of the determinant infimum produces arbitrarily good
F\o lner sets for every finite $S$.

\section{A quantitative compression lemma}

\begin{lemma}[Compression gap for the logarithm]\label{lem:compression}
Let $B$ be a bounded positive invertible operator on a Hilbert space, let
$P$ be a finite-rank orthogonal projection, and regard $PBP$ as an operator
on $P\mathcal H$.  Then
\begin{equation}\label{eq:compression-gap}
 \Tr_{P\mathcal H}\log(PBP)
 -\Tr_{P\mathcal H}(P\log(B)P)
 \geq \frac{\|(1-P)BP\|_{\HS}^2}{2\|B\|^2}.
\end{equation}
\end{lemma}

\begin{proof}
Put $Q=1-P$, and write $B$ in block form relative to
$\mathcal H=P\mathcal H\oplus Q\mathcal H$:
\[
 B=\begin{bmatrix}A&X^*\\X&C\end{bmatrix},
 \qquad A=PBP,\quad X=QBP,\quad C=QBQ.
\]
For $r\geq0$, the Schur complement formula gives
\[
 P(B+r)^{-1}P=(D_r-Y_r)^{-1},
 \quad D_r=A+rP,
 \quad Y_r=X^*(C+rQ)^{-1}X.
\]
Both $D_r$ and $D_r-Y_r$ are positive invertible.  If
$Z=D_r^{-1/2}Y_rD_r^{-1/2}$, then $0\leq Z<1$ and
$(1-Z)^{-1}\geq1+Z$.  Conjugating this inequality yields
\[
 (D_r-Y_r)^{-1}-D_r^{-1}\geq D_r^{-1}Y_rD_r^{-1}.
\]
Set $M=\|B\|$.  Since $C+rQ\leq(M+r)Q$ and
$D_r\leq(M+r)P$, we obtain
\begin{align*}
 \Tr\bigl(P(B+r)^{-1}P-(A+rP)^{-1}\bigr)
 &\geq \Tr(D_r^{-1}Y_rD_r^{-1})\\
 &\geq (M+r)^{-2}\Tr(Y_r)\\
 &\geq (M+r)^{-3}\|X\|_{\HS}^2.
\end{align*}
The resolvent representation
\[
 \log x=\int_0^\infty\left(\frac1{1+r}-\frac1{x+r}\right)\,dr
\]
now gives
\begin{align*}
 \Tr\log A-\Tr(P\log(B)P)
 &=\int_0^\infty
   \Tr\bigl(P(B+r)^{-1}P-(A+rP)^{-1}\bigr)\,dr\\
 &\geq \|X\|_{\HS}^2\int_0^\infty(M+r)^{-3}\,dr
 =\frac{\|X\|_{\HS}^2}{2M^2}.
\end{align*}
This is \eqref{eq:compression-gap}.
\end{proof}

\section{From the determinant gap to a group boundary}

For $s\in\G$, let $R_s\delta_g=\delta_{gs}$ be right translation on
$\ell^2(\G)$.  It commutes with $\lambda$.  Fix a finite symmetric
$S\subset\G\setminus\{e\}$, put
\[
 H=\sum_{s\in S}R_s,
 \qquad T=1+tH,
 \qquad 0<t<|S|^{-1},
\]
and set
\[
 \alpha=1-t|S|,
 \qquad \beta=1+t|S|.
\]
Since $H=H^*$ and $\|H\|\leq|S|$, we have
\begin{equation}\label{eq:T-bounds}
 0<\alpha\,1\leq T\leq\beta\,1.
\end{equation}
The group-ring element $a=1+t\sum_{s\in S}s$ is therefore positive and
invertible, and its right-convolution operator is $T$ (up to the harmless
inverse convention for right translation).

\begin{lemma}[Boundary lower bound]\label{lem:boundary}
Let $\phi=\tau(a^*(\,\cdot\,)a)$.  For every finite nonempty
$F\subset\G$,
\begin{equation}\label{eq:boundary-gap}
 \log\frac{(\det\phi[F])^{1/|F|}}{\Delta\phi}
 \geq
 \frac{2t^2(1-t|S|)^2}{(1+t|S|)^4}\,
 \frac1{|F|}\sum_{s\in S}|Fs\setminus F|.
\end{equation}
\end{lemma}

\begin{proof}
Let $P=P_F$ be the coordinate projection onto $\ell^2(F)$ and $Q=1-P$.
The vectors $\lambda(g)a\xi$, $g\in F$, are exactly $T\delta_g$, up to the
same inverse convention.  Thus their Gram matrix is
\begin{equation}\label{eq:gram-compression}
 \phi[F]=P T^2P\big|_{P\ell^2(\G)}
\end{equation}
up to simultaneous transposition, which does not change the determinant.
The source's definition gives
\begin{equation}\label{eq:delta-normalization}
 \log\Delta\phi=2\log\Delta T=\tau(\log T^2).
\end{equation}
Because $T^2$ commutes with $\lambda$, every diagonal coefficient of
$\log T^2$ is $\tau(\log T^2)$.  Hence
\[
 \Tr(P\log(T^2)P)=|F|\log\Delta\phi.
\]
Applying Lemma~\ref{lem:compression} to $B=T^2$ and using
$\|T^2\|^2\leq\beta^4$ yields
\begin{equation}\label{eq:first-gap}
 \log\frac{(\det\phi[F])^{1/|F|}}{\Delta\phi}
 \geq \frac{\|QT^2P\|_{\HS}^2}{2\beta^4|F|}.
\end{equation}

It remains to control the off-diagonal block.  Write
\[
 A_0=PTP,\qquad C_0=QTQ,\qquad Y=QTP.
\]
Then $A_0,C_0\geq\alpha$ by \eqref{eq:T-bounds}, and block multiplication
gives
\[
 QT^2P=C_0Y+YA_0.
\]
With the Hilbert--Schmidt inner product,
\begin{align*}
 2\alpha\|Y\|_{\HS}^2
 &\leq \Tr(Y^*C_0Y)+\Tr(Y^*YA_0)\\
 &=\langle Y,C_0Y+YA_0\rangle_{\HS}
 \leq \|Y\|_{\HS}\,\|QT^2P\|_{\HS}.
\end{align*}
Consequently
\begin{equation}\label{eq:sylvester}
 \|QT^2P\|_{\HS}\geq2\alpha\|QTP\|_{\HS}
 =2\alpha t\|QHP\|_{\HS}.
\end{equation}

For a fixed $g\in F$, the vectors $\delta_{gs}$, $s\in S$, are distinct.
After projecting outside $F$, they remain orthogonal.  Summing the squared
column norms gives the exact identity
\begin{equation}\label{eq:boundary-identity}
 \|QHP\|_{\HS}^2
 =\sum_{g\in F}\#\{s\in S:gs\notin F\}
 =\sum_{s\in S}|Fs\setminus F|.
\end{equation}
Combining \eqref{eq:first-gap}, \eqref{eq:sylvester}, and
\eqref{eq:boundary-identity} proves \eqref{eq:boundary-gap}.
\end{proof}

\section{Completion of the proof}

Assume the hypothesis of Theorem~\ref{thm:main}.  Fix an arbitrary finite
symmetric $S\subset\G\setminus\{e\}$ and any $0<t<|S|^{-1}$.  For the
corresponding $a$ and $\phi$, equality in \eqref{eq:source-lb} says that
there are finite nonempty sets $F_n$ for which
\[
 \log\frac{(\det\phi[F_n])^{1/|F_n|}}{\Delta\phi}\longrightarrow0.
\]
Lemma~\ref{lem:boundary} and positivity of its constant imply
\[
 \frac1{|F_n|}\sum_{s\in S}|F_ns\setminus F_n|\longrightarrow0.
\]
In particular $|F_ns\setminus F_n|/|F_n|\to0$ for every $s\in S$.
Since $S=S^{-1}$, this is equivalent to
$|F_n\triangle F_ns|/|F_n|\to0$ for every $s\in S$.

Every finite subset of $\G$ is contained in a finite symmetric one after
adjoining inverses.  We have therefore established the F\o lner criterion
for $\G$, so $\G$ is amenable.  This proves Theorem~\ref{thm:main}.

Equivalently, the contrapositive is quantitative.  If $\G$ is nonamenable,
then some finite symmetric $S$ has
\[
 \inf_{\varnothing\ne F\Subset\G}
 \frac1{|F|}\sum_{s\in S}|Fs\setminus F|>0.
\]
For every $0<t<|S|^{-1}$, Lemma~\ref{lem:boundary} supplies a uniform strict
gap in \eqref{eq:source-lb} for the single element
$a=1+t\sum_{s\in S}s$.  This also gives a concrete relation to the
isoperimetric profile requested in Problem 3.8 of the source.

\section{Verification, scope, and novelty}

The proof uses only positive invertible group-ring elements, so no kernel,
extended-logarithm, or singular-part convention is involved.  Equations
\eqref{eq:gram-compression} and \eqref{eq:delta-normalization} match the
source's Radon--Nikodym operator and determinant normalization: for a scalar
coefficient represented by $T\xi$, the source defines
$\Delta\phi=(\Delta T)^2$.

The script \texttt{code/compression\_gap\_probe.py} checks
Lemma~\ref{lem:compression} on 300 random complex positive matrices, the
Sylvester estimate \eqref{eq:sylvester} on another 300 random positive
matrices, and the boundary identity on a finite cyclic model without
wraparound.  These checks passed.  They are regression evidence, not part of
the proof.

The run's four cheap indexes and local source corpus were searched for the
arXiv id, title, and combinations of ``amenability'', ``finite compression'',
and ``Fuglede--Kadison determinant''.  Bounded exact-phrase, title/citation,
and arXiv web searches on July 21, 2026 found arXiv:2412.13751v4, last revised
March 20, 2026, where the question remains Problem 3.7, but no later paper
claiming a resolution.  This is therefore a candidate full solution with
moderate, not definitive, novelty confidence.

For human review, the decisive points are Lemma~\ref{lem:compression}, the
right-convolution identification in \eqref{eq:gram-compression}, and the
normalization in \eqref{eq:delta-normalization}.  The argument does not claim
a sharp constant in \eqref{eq:boundary-gap}, nor does it classify all
group-ring elements that witness nonamenability.

\begin{thebibliography}{9}
\bibitem{Austin}
T. Austin,
\emph{Entropy and determinants for unitary representations},
arXiv:2412.13751v4, 2026.
\end{thebibliography}

\end{document}

